% =====================================================================
%  Atividade 3 - Revisão bibliográfica: abordagens teóricas
%  Compilar no Overleaf com pdfLaTeX (padrão)
% =====================================================================
\documentclass[12pt,a4paper]{article}

\usepackage[utf8]{inputenc}
\usepackage[T1]{fontenc}
\usepackage[brazil]{babel}
\addto\captionsbrazil{\renewcommand{\tablename}{Quadro}}

\usepackage{mathptmx}
\usepackage[top=3cm, left=3cm, bottom=2cm, right=2cm]{geometry}
\usepackage{setspace}
\onehalfspacing          % 1,5 linha (ABNT). Se passar de 3 páginas: \setstretch{1.15}
\usepackage{indentfirst}
\setlength{\parindent}{1.25cm}
\usepackage{microtype}

\usepackage{booktabs}
\usepackage{tabularx}
\usepackage{titlesec}
\titleformat{\section}{\normalfont\bfseries}{\thesection}{0.5em}{}
\titleformat{\subsection}{\normalfont\bfseries\itshape}{\thesubsection}{0.5em}{}
\titlespacing*{\section}{0pt}{10pt}{4pt}
\titlespacing*{\subsection}{0pt}{6pt}{2pt}
\usepackage[hidelinks]{hyperref}

\newenvironment{referencias}{%
  \section*{Referências}
  \begin{singlespace}
  \setlength{\parindent}{0pt}
  \setlength{\parskip}{6pt}
  \raggedright
}{%
  \end{singlespace}
}

\begin{document}

\begin{center}
  {\bfseries ATIVIDADE 3 -- REVISÃO BIBLIOGRÁFICA: ABORDAGENS TEÓRICAS\\
  MEGAEVENTOS E EMPREGO EM ALOJAMENTO E ALIMENTAÇÃO NO PARÁ}\\[6pt]
  Marcela Paz da Silva Enriquez
\end{center}

\section{Introdução}

A COP30, realizada em Belém em novembro de 2025, foi um choque temporário e
concentrado de demanda sobre a economia paraense. Os setores de Alojamento e
Alimentação (divisões 55 e 56 da CNAE 2.0), que concentram o gasto de
visitantes, são tomados como \textit{proxy} do núcleo das atividades
turísticas. A pergunta da pesquisa é se, e em que magnitude, o fim do evento
esteve associado a uma redução do saldo mensal de empregos formais nesses
setores. Esta revisão apresenta as abordagens que explicam por que um evento
afetaria o emprego, por que o efeito se concentraria na hospitalidade e por
que seria temporário, e justifica o arcabouço adotado.

\section{Abordagens teóricas possíveis}

\subsection{Impacto econômico de megaeventos: promessas \textit{ex ante} e evidência \textit{ex post}}

Estudos \textit{ex ante}, baseados em matrizes insumo-produto e
multiplicadores, tendem a prever grandes ganhos; Crompton (1995) e Matheson
(2009) mostram as fontes dessa superestimação. No Brasil, Domingues,
Betarelli Junior e Magalhães (2011) projetaram alta de 0,5\% no emprego com a
Copa de 2014. Os estudos \textit{ex post} raramente confirmam esses números.
Hagn e Maennig (2009), comparando mês a mês o desemprego das 12 cidades-sede da
Copa de 2006 com o de 63 cidades alemãs, não encontram efeito significativo, e
observam que a cidade inteira pode ser uma unidade agregada demais para captar
o choque. Mattei e Mello (2017) não encontram diferenças relevantes entre
estados-sede e não sede da Copa de 2014. Os mecanismos recorrentes são a
substituição do gasto local, o \textit{crowding out} de turistas regulares e o
vazamento de renda para fora da região.

\subsection{Efeitos setoriais, temporários e dependentes do contexto}

Quando há efeito, ele é setorial e curto. Com dados mensais de 16 setores,
Feddersen e Maennig (2013) encontram aumento do emprego em Atlanta apenas no
mês das Olimpíadas de 1996, concentrado em comércio, alojamento e alimentação e
entretenimento. Boto-García e Santana-Gallego (2025) confirmam, com métodos
recentes, ganhos de emprego em serviços antes das Olimpíadas e Copas, que não
persistem após o evento. O caso mais próximo da COP30 é um evento não
esportivo: Bronzini, Mocetti e Mongardini (2019) mostram que o Jubileu de 2000
elevou a taxa de emprego de Roma, com o ganho concentrado em construção,
comércio e hotéis e restaurantes, sem efeito sobre o valor adicionado per
capita nem sobre o turismo no longo prazo. O contexto também importa: Wan e
Song (2019) encontram resultados distintos para Londres 2012 e para a Copa de
2014 e a Rio 2016, e Kobierecki e Pierzgalski (2022) não identificam efeito
sobre o crescimento das sedes entre 2010 e 2016.

\subsection{Demanda derivada e ajuste do emprego}

A demanda por trabalho deriva da demanda pelo serviço (HICKS, 1932): um hotel
contrata porque vende diárias. O megaevento é, assim, um choque temporário de
demanda; como as empresas contratam pela demanda esperada, as admissões
antecipam o evento, e os desligamentos seguem a demanda realizada. Alemayehu e
Tveteraas (2020) mostram, com dados diários de hotéis e restaurantes, que o
emprego se ajusta à demanda com defasagem e de forma incompleta, sobretudo em
alimentos e bebidas, o que sugere reações graduais, e não instantâneas, ao fim
do evento.

\subsection{Fluxos de emprego e rotatividade}

O emprego resulta de fluxos simultâneos de criação e destruição de postos,
muito maiores que a variação líquida (DAVIS; HALTIWANGER, 1992), que reagem a
choques de demanda (MORTENSEN; PISSARIDES, 1994). Dogru \textit{et al.} (2023)
mostram que a rotatividade na hospitalidade é a mais sensível ao ciclo da
economia. Isso justifica analisar admissões e desligamentos mensais e
controlar pela atividade econômica geral.

\subsection{Legado e economia política urbana}

Preuss (2007) analisa o legado de megaeventos em múltiplas dimensões, e Harvey
(1989) interpreta-os como estratégia de empreendedorismo urbano, com ganhos
concentrados e custos socializados. Essas abordagens enriquecem a
interpretação, mas não geram hipóteses testáveis com dados mensais de emprego.

\section{Escolha do arcabouço teórico}

Adota-se como arcabouço principal a literatura \textit{ex post} sobre
megaeventos (seções 2.1 e 2.2), microfundamentada pela demanda derivada (2.3)
e pela abordagem de fluxos (2.4); a perspectiva de legado (2.5) entra na
interpretação. A escolha se justifica pela adequação ao problema (efeito curto,
setorial e medido em fluxos de emprego formal), pela testabilidade (hipóteses
que se traduzem em um modelo com variáveis indicadoras para cada fase) e pelo
descarte fundamentado das alternativas: os estudos \textit{ex ante}
superestimam impactos, e a abordagem crítica opera na escala da governança
urbana. O Quadro~\ref{tab:arcabouco} resume essa correspondência.

\begin{table}[h]
  \centering
  \caption{Escolhas metodológicas e fundamentos teóricos}
  \label{tab:arcabouco}
  \small
  \begin{tabularx}{\textwidth}{Xl}
    \toprule
    \textbf{Escolha metodológica} & \textbf{Fundamento teórico} \\
    \midrule
    Fases: preparação, evento e desmobilização & Megaeventos \textit{ex post} \\
    Foco em Alojamento e Alimentação & Efeitos setoriais (FEDDERSEN; MAENNIG, 2013) \\
    Controle pelo volume de serviços & Demanda derivada; Dogru \textit{et al.} (2023) \\
    Saldo mensal de admissões e desligamentos & Fluxos de emprego e rotatividade \\
    Pará e Região Metropolitana de Belém & Agregação espacial (HAGN; MAENNIG, 2009) \\
    Comparação futura com outras capitais & Contrafactuais (BRONZINI \textit{et al.}, 2019) \\
    \bottomrule
  \end{tabularx}
\end{table}

\textbf{H1}: o saldo de empregos aumenta imediatamente antes e durante o
evento. \textbf{H2}: controlados sazonalidade e volume de serviços, o saldo na
desmobilização fica abaixo do período de referência. \textbf{H3}: os efeitos se
concentram na Região Metropolitana de Belém e no alojamento.

\section{Implicações para a base de dados e para as inferências}

A base usa os microdados do Novo CAGED (admissões e desligamentos formais por
subclasse CNAE e município, pelo mês de competência; BRASIL, 2023) e o índice
de volume de serviços do Pará (IBGE, PMS), de janeiro de 2021 a julho de 2026.
A série começa após a quebra do eSocial em 2020 (IPEA, 2020). Ela permite
inferir se o saldo se desviou do padrão esperado em cada fase e onde o desvio
foi maior, mas, sem grupo de controle, a inferência é de associação
condicional. A literatura aponta o passo seguinte: construir um contrafactual
com outras capitais, por diferenças em diferenças ou controle sintético, como
em Hagn e Maennig (2009) e Bronzini, Mocetti e Mongardini (2019).

\section{Lacuna e contribuição}

Estudos quantitativos sobre eventos não esportivos são raros (BRONZINI;
MOCETTI; MONGARDINI, 2019), e em levantamento preliminar não foram
identificados estudos \textit{ex post} sobre o emprego local em conferências
climáticas. A pesquisa contribui ao analisar a COP30 em uma cidade amazônica,
com microdados mensais do emprego formal em um país em desenvolvimento.

\begin{referencias}

ALEMAYEHU, F. K.; TVETERAAS, S. L. Long-run labour flexibility in
hospitality: a dynamic common correlated effects approach. \textbf{Tourism
Economics}, v. 26, n. 4, p. 704-718, 2020.

BOTO-GARCÍA, D.; SANTANA-GALLEGO, M. Old questions, new methods: revisiting
the economic effects of hosting mega-sport events. \textbf{Tourism
Management}, v. 113, 105324, 2025.

BRASIL. Ministério do Trabalho e Emprego. \textbf{Portaria MTE nº 2.420, de
10 de julho de 2023}. Anexo XV: metodologia do Novo CAGED. Brasília, 2023.

BRONZINI, R.; MOCETTI, S.; MONGARDINI, M. \textbf{The economic effects of
big events}: evidence from the Great Jubilee 2000 in Rome. Roma: Banca
d'Italia, 2019. (Temi di discussione, n. 1208).

CROMPTON, J. L. Economic impact analysis of sports facilities and events:
eleven sources of misapplication. \textbf{Journal of Sport Management}, v. 9,
n. 1, p. 14-35, 1995.

DAVIS, S. J.; HALTIWANGER, J. Gross job creation, gross job destruction, and
employment reallocation. \textbf{The Quarterly Journal of Economics}, v. 107,
n. 3, p. 819-863, 1992.

DOGRU, T.; McGINLEY, S.; SHARMA, A.; IŞIK, C.; HANKS, L. Employee turnover
dynamics in the hospitality industry vs. the overall economy.
\textbf{Tourism Management}, v. 99, 104783, 2023.

DOMINGUES, E. P.; BETARELLI JUNIOR, A. A.; MAGALHÃES, A. S. Quanto vale o
show? Impactos econômicos dos investimentos da Copa do Mundo 2014 no Brasil.
\textbf{Estudos Econômicos}, v. 41, n. 2, p. 409-439, 2011.

FEDDERSEN, A.; MAENNIG, W. Mega-events and sectoral employment: the case of
the 1996 Olympic Games. \textbf{Contemporary Economic Policy}, v. 31, n. 3,
p. 580-603, 2013.

HAGN, F.; MAENNIG, W. Large sport events and unemployment: the case of the
2006 soccer World Cup in Germany. \textbf{Applied Economics}, v. 41, n. 25,
p. 3295-3302, 2009.

HARVEY, D. From managerialism to entrepreneurialism: the transformation in
urban governance in late capitalism. \textbf{Geografiska Annaler: Series B,
Human Geography}, v. 71, n. 1, p. 3-17, 1989.

HICKS, J. R. \textbf{The theory of wages}. London: Macmillan, 1932.

IPEA. Substituição da captação dos dados do Caged pelo eSocial: implicações
para as estatísticas do emprego formal. \textbf{Boletim Mercado de Trabalho:
conjuntura e análise}, n. 69, 2020.

KOBIERECKI, M. M.; PIERZGALSKI, M. Sports mega-events and economic growth: a
synthetic control approach. \textbf{Journal of Sports Economics}, v. 23,
n. 5, p. 567-597, 2022.

MATHESON, V. A. Economic multipliers and mega-event analysis.
\textbf{International Journal of Sport Finance}, v. 4, n. 1, p. 63-70, 2009.

MATTEI, T. S.; MELLO, G. R. Copa do mundo de futebol 2014: análise do
impacto da receita e volume de vendas do comércio, dos serviços e da produção
industrial nos estados brasileiros. \textbf{Revista Gestão \& Tecnologia},
v. 17, n. 1, p. 207-231, 2017.

MORTENSEN, D. T.; PISSARIDES, C. A. Job creation and job destruction in the
theory of unemployment. \textbf{The Review of Economic Studies}, v. 61, n. 3,
p. 397-415, 1994.

PREUSS, H. The conceptualisation and measurement of mega sport event
legacies. \textbf{Journal of Sport \& Tourism}, v. 12, n. 3-4, p. 207-228,
2007.

WAN, S. K.; SONG, H. Economic impact assessment of mega-events in the United
Kingdom and Brazil. \textbf{Journal of Hospitality \& Tourism Research},
v. 43, n. 7, p. 1044-1067, 2019.

\end{referencias}

\end{document}
